\section{Design}

This section presents the architecture of \axvisor{} as an OS-portable
hypervisor. We first state the design goals and the principles that govern the
substrate contract (\S\ref{sec:design-goals}). We then describe the overall
architecture (\S\ref{sec:arch-overview}), the substrate interface contract in
detail (\S\ref{sec:contract}), and the guest-facing KVM-compatible ABI that
completes the decoupling (\S\ref{sec:kvm-abi}).

\subsection{Design Goals and Principles}
\label{sec:design-goals}

\axvisor{} is guided by four goals, each grounded in the coupling problem
identified in Section~\ref{sec:design-goals}.

\paragraph{G1. Substrate-OS portability}
The same hypervisor core must run on substrates that differ in kernel
architecture, memory model, scheduling model, and device infrastructure.
Portability is measured by how little substrate-specific code is required, not
merely by the fact that the system can be made to compile on each target.

\paragraph{G2. Minimal, faithful contract}
The dependency of the hypervisor core on its substrate must be expressed as a
narrow, explicit interface. The interface must expose services at the
granularity virtualization requires---physical memory, execution contexts,
interrupts---rather than at the granularity a particular OS happens to offer. It
must not hide semantic differences that affect correctness.

\paragraph{G3. Guest-side neutrality}
Guests and their toolchains must not be coupled to the substrate choice. By
exposing a KVM-compatible ABI, \axvisor{} lets unmodified guests and existing
VMMs run on top of any substrate that implements the contract.

\paragraph{G4. Bounded adaptation cost}
Implementing the contract on a new substrate should require proportionally less
effort than writing a hypervisor from scratch. The adaptation layer, not the
core, should absorb substrate idiosyncrasies.

Three principles follow. First, \emph{separation of virtualization logic from OS
services}: anything that is genuinely about virtualization---vCPU state, guest
memory bookkeeping, trap dispatch, virtual device mediation---stays in the core;
anything that is about obtaining a resource or invoking a platform policy lives
behind the contract. Second, \emph{explicit over implicit}: where substrates
differ in semantics, the contract surfaces the difference rather than silently
normalizing it, so that adaptation layers make conscious choices. Third,
\emph{contract by capability}: a substrate is eligible to host \axvisor{} if and
only if it can provide the contract's required capabilities; optional
capabilities may be absent without breaking the core, but they bound the
hypervisor's feature set on that substrate.

\subsection{Architecture Overview}
\label{sec:arch-overview}

\axvisor{} is structured into three layers, shown in Figure~\ref{fig:arch}. At
the top, one or more guest VMs interact with the hypervisor through the
KVM-compatible ABI. In the middle, the portable hypervisor core implements all
virtualization logic that is independent of any substrate: vCPU lifecycle and
state machine, stage-2/EPT guest memory management, trap and exception dispatch,
hypercall handling, virtual device models, and VM lifecycle. At the bottom, a
substrate adaptation layer implements the substrate interface contract for a
specific OS. The core calls only into the contract; the adaptation layer calls
into the substrate's native APIs.

\begin{figure}[t]
  \centering
  \framebox[0.95\columnwidth]{\parbox{0.9\columnwidth}{\centering
    \vspace{0.5em}
    Guest VMs\\ $\updownarrow$\\ KVM ABI\\ $\updownarrow$\\
    Portable Core\\ $\updownarrow$\\ Substrate Interface Contract\\ $\updownarrow$\\
    \{\arceos{}, \asterinas{}, Linux\} adaptation layers\\ $\updownarrow$\\
    Hardware
    \vspace{0.5em}}}
  \caption{The \axvisor{} architecture. The portable core depends only on the
  substrate interface contract; each adaptation layer implements the contract
  for one OS. Guests interact through a KVM-compatible ABI, decoupling them from
  the substrate choice.}
  \label{fig:arch}
\end{figure}

This layering converts the substrate OS from an implicit, ambient dependency
into an explicit service provider. The boundary between core and adaptation
layer is the single seam along which portability is won or lost; everything
above it is shared, everything below it is substrate-specific. Crucially, the
KVM-compatible ABI provides a second, independent seam on the guest side, so
that guests and substrates can vary independently of each other.

\subsection{The Substrate Interface Contract}
\label{sec:contract}

The contract is the heart of the design. We define it as a set of service
groups, each containing a small number of operations. The contract is
deliberately not a POSIX-like API: hypervisors need capabilities that
general-purpose application interfaces neither expose nor guarantee, such as
allocating contiguous physical memory, installing stage-2 mappings, creating
privileged execution contexts, and registering direct interrupt handlers. We
therefore call these \emph{hypervisor-grade services}.

\paragraph{Memory management}
The core must allocate physical memory regions of defined size and contiguity,
map guest physical addresses to host physical frames in the stage-2 page table,
and manage ownership and lifetime of frames handed to guests. The contract
exposes: \code{alloc\_phys} (contiguous or discontiguous physical frames),
\code{map\_stage2} (install a guest-physical to host-physical mapping with
permissions), \code{unmap\_stage2}, and \code{flush\_tlb}. Semantic pressure
here is high: substrates differ in whether they expose physical frame numbers
directly, whether they permit fine-grained stage-2 control, and how they handle
cache coherence and DMA ownership. The contract requires the substrate to
provide physical addresses and stage-2 control; if a substrate abstracts these
away, its adaptation layer must recover them.

\paragraph{CPU and vCPU execution}
The core creates vCPUs as execution contexts and drives them through
run/pause/halt cycles, reading and writing guest register state. The contract
exposes: \code{vcpu\_create}, \code{vcpu\_run} (enter the guest until a trap),
\code{vcpu\_get\_regs}/\code{vcpu\_set\_regs}, and \code{vcpu\_destroy}. The
substrate must provide a way to enter a hardware virtualization mode (e.g., EL2
on Arm, VMX root on x86) and to surface traps back to the core. How the
substrate schedules the host thread that backs a vCPU is a substrate policy; the
contract only requires that \code{vcpu\_run} blocks until the guest traps.

\paragraph{Interrupts and timers}
The core must inject virtual interrupts into guests, receive physical interrupts
that target assigned devices, and arm timers for scheduling quantum and watchdog
purposes. The contract exposes: \code{irq\_register} (bind a physical interrupt
to a core handler), \code{irq\_inject} (deliver a virtual interrupt to a vCPU),
\code{irq\_eoi} (end-of-interrupt), and \code{timer\_arm}. Interrupt handling is
among the most semantically divergent areas: substrates differ in
interrupt-controller models, in whether handlers run in interrupt context or
deferred context, and in how interrupt affinity is configured. The contract
requires the substrate to deliver a physical interrupt to a core-registered
callback and to honor injection requests; the adaptation layer reconciles
controller specifics.

\paragraph{Synchronization}
The core uses locks and barriers to protect shared data structures across vCPUs
and interrupt handlers. The contract exposes spinlocks and a minimal
memory-barrier facility. Because some substrates disable preemption or run with
interrupts masked in core paths, the contract requires that spinlocks be usable
in interrupt context and that they provide ordering compatible with the target
ISA's memory model.

\paragraph{Device and MMIO access}
The core accesses assigned devices via MMIO and mediates emulated devices for
guests. The contract exposes: \code{mmio\_read}/\code{mmio\_write} (physical
MMIO access), \code{io\_port} (x86 I/O space, where applicable), and
\code{register\_mmio\_handler} (install a core handler for a guest MMIO region
that triggers on guest access). For passthrough, the substrate must permit
direct device assignment; for emulation, the substrate must trap guest MMIO
faults and forward them.

\paragraph{DMA and IOMMU}
For device passthrough and shared-memory communication, the core must program
IOMMU mappings so that assigned devices can access guest memory safely. The
contract exposes: \code{iommu\_map}/\code{iommu\_unmap} and
\code{iommu\_attach\_device}. Where a substrate lacks an IOMMU or does not
expose it, this capability is optional and passthrough is constrained
accordingly; the core degrades gracefully rather than failing.

\paragraph{Lifecycle and logging}
The contract provides initialization entry points, a way to query the
platform's CPU/memory layout at boot, and logging primitives. These are
necessary but low-risk; their main purpose is to give the adaptation layer a
defined place to hand control to the core.

Table~\ref{tab:contract} summarizes the contract. The key property is that every
operation is expressed in terms of what virtualization needs, not how a
particular OS provides it. A substrate that can supply these
capabilities---regardless of whether it is a unikernel, a monolithic kernel, or
a framekernel---is eligible to host the core.

\begin{table}[t]
  \centering
  \caption{Substrate interface contract: service groups and representative
  operations.}
  \label{tab:contract}
  \begin{tabular}{@{}p{2.4cm}p{5.0cm}@{}}
    \toprule
    \textbf{Service group} & \textbf{Representative operations} \\
    \midrule
    Memory & \code{alloc\_phys}, \code{map\_stage2}, \code{unmap\_stage2},
      \code{flush\_tlb} \\
    CPU/vCPU & \code{vcpu\_create}, \code{vcpu\_run},
      \code{vcpu\_get/set\_regs}, \code{vcpu\_destroy} \\
    Interrupt/timer & \code{irq\_register}, \code{irq\_inject},
      \code{irq\_eoi}, \code{timer\_arm} \\
    Synchronization & \code{spinlock}, memory barriers \\
    Device/MMIO & \code{mmio\_read/write}, \code{register\_mmio\_handler} \\
    DMA/IOMMU & \code{iommu\_map/unmap}, \code{iommu\_attach\_device}
      (optional) \\
    Lifecycle & \code{init}, \code{query\_platform}, \code{log} \\
    \bottomrule
  \end{tabular}
\end{table}

\subsection{KVM-Compatible ABI}
\label{sec:kvm-abi}

Decoupling the substrate side is necessary but not sufficient: if guests and
VMMs had to be rewritten for each substrate, portability would be hollow.
\axvisor{} therefore exposes a KVM-compatible ABI on the guest-facing side. The
KVM ioctl interface---operations such as \code{KVM\_CREATE\_VM},
\code{KVM\_CREATE\_VCPU}, \code{KVM\_RUN}, and the memory-slot and irqchip
families---has become the de facto standard for VM
management~\cite{kivity2007kvm}. A large body of guest drivers, virtio backends,
and management toolchains assume it.

\axvisor{} implements a KVM-ABI shim inside the portable core. The shim
translates standard KVM ioctls into internal core operations that are themselves
substrate-independent. When a VMM issues \code{KVM\_CREATE\_VCPU}, the core
invokes \code{vcpu\_create} on the contract; when it issues \code{KVM\_RUN}, the
core invokes \code{vcpu\_run}. Because both the ioctl surface and the core logic
are substrate-neutral, the same guest and the same VMM run on top of \arceos{},
\asterinas{}, or Linux without modification. The KVM ABI thus acts as the
guest-side counterpart to the substrate contract: just as the contract decouples
the core from the OS below, the ABI decouples guests from the substrate choice
below the core.

There is a deliberate asymmetry worth making explicit. On the substrate side,
the contract is a \emph{provider} interface: each OS implements it. On the guest
side, the KVM ABI is a \emph{consumer} interface: the core implements it for
guests to consume. The portable core sits between two contracts, neither of
which references a specific OS. This double decoupling is what allows the matrix
of guests $\times$ substrates to be supported without an $O(g \times s)$
implementation effort: the core is written once, each substrate contributes one
adaptation layer, and all guests are inherited from the existing KVM ecosystem.
